\documentclass{article}
\usepackage{amsmath,amsfonts,amssymb}
\usepackage{algorithm}
\usepackage{algpseudocode}
\usepackage{hyperref}
\usepackage{bm}
\usepackage{enumitem}
\usepackage{graphicx}
\usepackage{color}
\usepackage{tikz}
\usepackage{booktabs}
\usepackage{array}
\usepackage{float}
\usepackage{caption}
\usepackage{fancyhdr}
\usepackage{geometry}
\geometry{margin=1in}
\pagestyle{plain}
\setlength{\parskip}{0.5em}
\setlength{\parindent}{0pt}
\begin{document}

%%%%%%%%%%%%%%%%%%%%%%%%%%%%%%%%%%%%%%%%%%%%%%%%%%%%%%%%%%%%
\section*{DP‑SGD with Gaussian Noise}
%%%%%%%%%%%%%%%%%%%%%%%%%%%%%%%%%%%%%%%%%%%%%%%%%%%%%%%%%%%%

\subsection*{Mathematical Formulation}
Let $f:\mathbb{R}^d\!\to\!\mathbb{R}$ be the (possibly non‑convex) loss
\[
    f(\mathbf{x}) \;=\; \frac{1}{n}\sum_{i=1}^{n}\ell(\mathbf{x};\mathbf{z}_i),
\]
where $\mathbf{z}_i$ denotes a data point.  
At iteration $t$ we sample a minibatch $\mathcal{B}_t\subset\{1,\dots,n\}$ of size $B$ uniformly at random and compute the *clipped* stochastic gradient
\[
    \tilde{\mathbf{g}}_t
    \;=\;\frac{1}{B}\sum_{i\in\mathcal{B}_t}
    \operatorname{clip}_C\!\bigl(\nabla\ell(\mathbf{x}_t;\mathbf{z}_i)\bigr),
    \qquad 
    \operatorname{clip}_C(\mathbf{u})\;=\;
    \mathbf{u}\,\min\!\Bigl\{1,\frac{C}{\|\mathbf{u}\|_2}\Bigr\},
\]
so that $\|\tilde{\mathbf{g}}_t\|_2\le C$ almost surely.  
Gaussian noise is added:
\[
    \mathbf{g}_t \;=\; \tilde{\mathbf{g}}_t \;+\; \mathbf{n}_t,
    \qquad
    \mathbf{n}_t\sim\mathcal{N}\bigl(\mathbf{0},\sigma^2\mathbf{I}_d\bigr).
\]
The parameter update is
\[
    \boxed{\;\mathbf{x}_{t+1}
    \;=\;\mathbf{x}_t \;-\; \eta_t \,\mathbf{g}_t\;}
\tag{1}
\]
with learning‑rate schedule $\{\eta_t\}_{t\ge0}$.
The sequence $\{(\eta_t,\sigma,C,B)\}$ together with the number of
iterations $T$ determines the differential‑privacy budget $(\varepsilon,\delta)$ via the moments accountant or Rényi DP; we treat $(\varepsilon,\delta)$ as a *post‑hoc* bookkeeping quantity and do not need its exact formula for the convergence analysis.

\subsection*{Pseudocode}
\begin{algorithm}[H]
\caption{DP‑SGD with Gaussian Noise}
\begin{algorithmic}[1]
\State \textbf{Input:} initial point $\mathbf{x}_0$, learning‑rate schedule $\{\eta_t\}$,
clipping bound $C$, noise scale $\sigma$, minibatch size $B$, total iterations $T$.
\For{$t=0,\dots,T-1$}
    \State Sample a minibatch $\mathcal{B}_t$ of size $B$ uniformly.
    \State Compute per‑example gradients $\mathbf{g}_i=\nabla\ell(\mathbf{x}_t;\mathbf{z}_i)$ for $i\in\mathcal{B}_t$.
    \State Clip: $\mathbf{c}_i=\operatorname{clip}_C(\mathbf{g}_i)$.
    \State Form the average $\tilde{\mathbf{g}}_t=\frac{1}{B}\sum_{i\in\mathcal{B}_t}\mathbf{c}_i$.
    \State Sample noise $\mathbf{n}_t\sim\mathcal{N}(\mathbf{0},\sigma^2\mathbf{I}_d)$.
    \State Noisy gradient $\mathbf{g}_t\gets \tilde{\mathbf{g}}_t+\mathbf{n}_t$.
    \State Update $\mathbf{x}_{t+1}\gets \mathbf{x}_t-\eta_t\mathbf{g}_t$ \Comment{Eq.~(1)}.
\EndFor
\State \textbf{Output:} $\mathbf{x}_T$ (or the average $\bar{\mathbf{x}}_T=\frac1T\sum_{t=0}^{T-1}\mathbf{x}_t$).
\end{algorithmic}
\end{algorithm}

\subsection*{Dry‑Run Example}
Consider the 1‑dimensional quadratic
\[
    f(w)=(w-3)^2,
\]
with exact gradient $\nabla f(w)=2(w-3)$.  
Take $w_0=0$, learning‑rate $\eta=0.1$, clipping bound $C=5$ (so clipping never activates), noise variance $\sigma^2=0$ (to keep the arithmetic simple), and a batch size $B=1$ (single deterministic gradient).  

\begin{center}
\begin{tabular}{c|c|c|c}
\textbf{Iter.} $t$ & $\nabla f(w_t)$ & $g_t=\nabla f(w_t)$ & $w_{t+1}=w_t-\eta g_t$ \\ \midrule
0 & $2(0-3)=-6$ & $-6$ & $0-0.1(-6)=0.6$ \\
1 & $2(0.6-3)=-4.8$ & $-4.8$ & $0.6-0.1(-4.8)=1.08$ \\
2 & $2(1.08-3)=-3.84$ & $-3.84$ & $1.08-0.1(-3.84)=1.464$ 
\end{tabular}
\end{center}
Corresponding objective values:
\[
f(0)=9,\quad f(0.6)=5.76,\quad f(1.08)=3.6864,\quad f(1.464)=2.350\;.
\]
The iterates move toward the minimiser $w^\star=3$, illustrating the update (1) in the absence of noise.

%%%%%%%%%%%%%%%%%%%%%%%%%%%%%%%%%%%%%%%%%%%%%%%%%%%%%%%%%%%%
\section*{Assumptions}
%%%%%%%%%%%%%%%%%%%%%%%%%%%%%%%%%%%%%%%%%%%%%%%%%%%%%%%%%%%%
\begin{enumerate}[label={A\arabic*}:, leftmargin=*, itemsep=0.3em]
\item \textbf{Smoothness.} $f$ is $L$‑smooth:
\[
    \forall\mathbf{x},\mathbf{y}\in\mathbb{R}^d,\qquad
    f(\mathbf{y})\le f(\mathbf{x})+\langle\nabla f(\mathbf{x}),\mathbf{y}-\mathbf{x}\rangle
    +\frac{L}{2}\|\mathbf{y}-\mathbf{x}\|_2^{2}.
\tag{A1}
\]
\item \textbf{Bounded gradient norm after clipping.}
\[
    \|\tilde{\mathbf{g}}_t\|_2\le C \quad\text{a.s.}
\tag{A2}
\]
\item \textbf{Unbiasedness of the clipped minibatch gradient.}
Conditioned on $\mathbf{x}_t$,
\[
    \mathbb{E}\!\bigl[\tilde{\mathbf{g}}_t\mid\mathbf{x}_t\bigr]=\nabla f(\mathbf{x}_t).
\tag{A3}
\]
(The expectation is over the random minibatch; clipping does not introduce bias when the per‑example gradients are already bounded by $C$, which we assume for simplicity.)
\item \textbf{Gaussian noise.}
\[
    \mathbb{E}[\mathbf{n}_t]=\mathbf{0},\qquad
    \mathbb{E}\bigl[\mathbf{n}_t\mathbf{n}_t^{\!\top}\bigr]=\sigma^{2}\mathbf{I}_d.
\tag{A4}
\]
\item \textbf{Learning‑rate schedule.} We use a non‑increasing step size
\[
    \eta_t=\frac{\eta}{\sqrt{t+1}},\qquad \eta>0.
\tag{A5}
\]
\item \textbf{Finite initial suboptimality.}
\[
    \Delta_0:=f(\mathbf{x}_0)-f^\star<\infty,\qquad f^\star:=\inf_{\mathbf{x}} f(\mathbf{x}).
\tag{A6}
\]
\end{enumerate}

%%%%%%%%%%%%%%%%%%%%%%%%%%%%%%%%%%%%%%%%%%%%%%%%%%%%%%%%%%%%
\section*{Main Convergence Theorem}
%%%%%%%%%%%%%%%%%%%%%%%%%%%%%%%%%%%%%%%%%%%%%%%%%%%%%%%%%%%%
\begin{theorem}[Expected Gradient Norm for DP‑SGD]\label{thm:main}
Under assumptions A1–A6, let $\{ \mathbf{x}_t\}_{t=0}^{T-1}$ be generated by
Algorithm~1 with the step size $\eta_t=\eta/\sqrt{t+1}$.
Define the averaged iterate $\bar{\mathbf{x}}_T:=\frac{1}{T}\sum_{t=0}^{T-1}\mathbf{x}_t$.
Then
\[
\frac{1}{T}\sum_{t=0}^{T-1}\mathbb{E}\!\bigl[\|\nabla f(\mathbf{x}_t)\|_2^{2}\bigr]
\;\le\;
\underbrace{\frac{2\Delta_0}{\eta\sqrt{T}}}_{\text{optimization term}}
\;+\;
\underbrace{\frac{L\eta C^{2}}{B}\,\sigma^{2}}_{\text{privacy‑noise term}}
\;+\;
\underbrace{\frac{L\eta^{2}C^{2}}{2}}_{\text{discretisation term}} .
\tag{2}
\]
Consequently, choosing $\eta = \Theta\!\bigl(T^{-1/4}\bigr)$ yields
\[
    \frac{1}{T}\sum_{t=0}^{T-1}\mathbb{E}\!\bigl[\|\nabla f(\mathbf{x}_t)\|_2^{2}\bigr]
    =\mathcal{O}\!\bigl(T^{-1/2}\bigr).
\]
\end{theorem}

\subsection*{Theoretical Properties}
\begin{itemize}[leftmargin=*]
\item \textbf{Stability.} The bound (2) shows that the additive Gaussian noise only contributes a *constant* term proportional to $\sigma^{2}$; decreasing $\eta$ mitigates its effect.
\item \textbf{Hyper‑parameter sensitivity.}
    \begin{itemize}
        \item Larger clipping $C$ improves the unbiasedness of the stochastic gradient (A3) but inflates the noise term $L\eta C^{2}\sigma^{2}/B$.
        \item Batch size $B$ appears only in the denominator of the noise term, confirming the well‑known variance‑reduction effect of larger minibatches.
        \item The step‑size $\eta$ balances three terms: a $1/\eta\sqrt{T}$ optimisation term, a $\eta$‑scaled noise term, and an $\eta^{2}$ discretisation term.
    \end{itemize}
\item \textbf{Comparison to non‑private SGD.}
    Setting $\sigma=0$ recovers the classical non‑convex SGD bound
    $\frac{2\Delta_0}{\eta\sqrt{T}}+\frac{L\eta^{2}C^{2}}{2}$,
    which is optimal up to constants for smooth non‑convex objectives.
\end{itemize}

%%%%%%%%%%%%%%%%%%%%%%%%%%%%%%%%%%%%%%%%%%%%%%%%%%%%%%%%%%%%
\section*{Preliminaries (Definitions and Lemmas)}
%%%%%%%%%%%%%%%%%%%%%%%%%%%%%%%%%%%%%%%%%%%%%%%%%%%%%%%%%%%%
\begin{lemma}[Smoothness inequality]\label{lem:smooth}
If $f$ satisfies A1, then for any $\mathbf{x},\mathbf{y}\in\mathbb{R}^{d}$,
\[
    f(\mathbf{y})\le f(\mathbf{x})+\langle\nabla f(\mathbf{x}),\mathbf{y}-\mathbf{x}\rangle
    +\frac{L}{2}\|\mathbf{y}-\mathbf{x}\|_2^{2}.
\tag{L1}
\]
\end{lemma}
\begin{proof}
Standard consequence of $L$‑smoothness; see e.g. Nesterov (2004). \hfill $\square$
\end{proof}

\begin{lemma}[Noise variance bound]\label{lem:noise}
For the Gaussian noise $\mathbf{n}_t$ defined in A4,
\[
    \mathbb{E}\!\bigl[\|\mathbf{n}_t\|_2^{2}\bigr]=d\sigma^{2}.
\tag{L2}
\]
If we regard the *per‑coordinate* variance, the bound $ \mathbb{E}[\|\mathbf{n}_t\|_2^{2}] \le \sigma^{2} $ holds after scaling $C$ into the definition of $\sigma$ (common DP literature). For brevity we absorb $d$ into $\sigma^{2}$, i.e. we treat $\sigma^{2}$ as the *effective* variance of the vector norm.
\end{lemma}
\begin{proof}
Direct computation for a zero‑mean isotropic Gaussian. \hfill $\square$
\end{proof}

%%%%%%%%%%%%%%%%%%%%%%%%%%%%%%%%%%%%%%%%%%%%%%%%%%%%%%%%%%%%
\section*{Main Proof (Step‑by‑Step Derivation)}
%%%%%%%%%%%%%%%%%%%%%%%%%%%%%%%%%%%%%%%%%%%%%%%%%%%%%%%%%%%%
We prove Theorem~\ref{thm:main}. Throughout the proof, expectations are taken over all randomness up to iteration $t$ (minibatch sampling and Gaussian noise).  

\noindent\textbf{Notation.} Write $\mathbf{g}_t = \tilde{\mathbf{g}}_t+\mathbf{n}_t$ for the noisy gradient and $\Delta_t:=f(\mathbf{x}_t)-f^\star$.

\paragraph{Step 1 – Apply smoothness.}
Using Lemma~\ref{lem:smooth} with $\mathbf{x}=\mathbf{x}_t$ and $\mathbf{y}=\mathbf{x}_{t+1}$,
\[
\begin{aligned}
f(\mathbf{x}_{t+1})
&\le f(\mathbf{x}_t)
   +\big\langle\nabla f(\mathbf{x}_t),\mathbf{x}_{t+1}-\mathbf{x}_t\big\rangle
   +\frac{L}{2}\|\mathbf{x}_{t+1}-\mathbf{x}_t\|_2^{2}
\\[4pt]
&= f(\mathbf{x}_t)
   -\eta_t\big\langle\nabla f(\mathbf{x}_t),\mathbf{g}_t\big\rangle
   +\frac{L}{2}\eta_t^{2}\|\mathbf{g}_t\|_2^{2}.
\tag{3}
\end{aligned}
\]
[Step 1]

\paragraph{Step 2 – Expand the inner product.}
Decompose $\mathbf{g}_t = \tilde{\mathbf{g}}_t+\mathbf{n}_t$ and use linearity:
\[
\big\langle\nabla f(\mathbf{x}_t),\mathbf{g}_t\big\rangle
 = \big\langle\nabla f(\mathbf{x}_t),\tilde{\mathbf{g}}_t\big\rangle
   +\big\langle\nabla f(\mathbf{x}_t),\mathbf{n}_t\big\rangle .
\tag{4}
\]
[Step 2]

\paragraph{Step 3 – Take conditional expectation.}
Condition on $\mathbf{x}_t$ and the minibatch; the noise is independent and zero‑mean (A4). Therefore
\[
\mathbb{E}\!\bigl[\langle\nabla f(\mathbf{x}_t),\mathbf{n}_t\rangle\mid\mathbf{x}_t\bigr]=0.
\tag{5}
\]
Moreover, by unbiasedness A3,
\[
\mathbb{E}\!\bigl[\tilde{\mathbf{g}}_t\mid\mathbf{x}_t\bigr]=\nabla f(\mathbf{x}_t)
\;\Longrightarrow\;
\mathbb{E}\!\bigl[\langle\nabla f(\mathbf{x}_t),\tilde{\mathbf{g}}_t\rangle\mid\mathbf{x}_t\bigr]
= \|\nabla f(\mathbf{x}_t)\|_2^{2}.
\tag{6}
\]
[Step 3]

\paragraph{Step 4 – Expectation of the squared norm.}
Using $\|\mathbf{a}+\mathbf{b}\|^{2}\le 2\|\mathbf{a}\|^{2}+2\|\mathbf{b}\|^{2}$,
\[
\|\mathbf{g}_t\|_2^{2}
 = \|\tilde{\mathbf{g}}_t+\mathbf{n}_t\|_2^{2}
 \le 2\|\tilde{\mathbf{g}}_t\|_2^{2}+2\|\mathbf{n}_t\|_2^{2}.
\tag{7}
\]
Taking expectation conditioned on $\mathbf{x}_t$ and using A2 ($\|\tilde{\mathbf{g}}_t\|\le C$) together with Lemma~\ref{lem:noise},
\[
\mathbb{E}\!\bigl[\|\mathbf{g}_t\|_2^{2}\mid\mathbf{x}_t\bigr]
\;\le\; 2C^{2}+2\sigma^{2}.
\tag{8}
\]
[Step 4]

\paragraph{Step 5 – Plug expectations into (3).}
Take total expectation of (3) and substitute (5)–(8):
\[
\begin{aligned}
\mathbb{E}\!\bigl[f(\mathbf{x}_{t+1})\bigr]
&\le \mathbb{E}\!\bigl[f(\mathbf{x}_t)\bigr]
   -\eta_t\,\mathbb{E}\!\bigl[\|\nabla f(\mathbf{x}_t)\|_2^{2}\bigr]
   +\frac{L}{2}\eta_t^{2}\,\mathbb{E}\!\bigl[\|\mathbf{g}_t\|_2^{2}\bigr]
\\[4pt]
&\le \mathbb{E}\!\bigl[f(\mathbf{x}_t)\bigr]
   -\eta_t\,\mathbb{E}\!\bigl[\|\nabla f(\mathbf{x}_t)\|_2^{2}\bigr]
   +\frac{L}{2}\eta_t^{2}\bigl(2C^{2}+2\sigma^{2}\bigr).
\end{aligned}
\tag{9}
\]
[Step 5]

\paragraph{Step 6 – Rearrange to isolate the gradient term.}
Move the gradient term to the left and sum from $t=0$ to $T-1$:
\[
\sum_{t=0}^{T-1}\eta_t\,
\mathbb{E}\!\bigl[\|\nabla f(\mathbf{x}_t)\|_2^{2}\bigr]
\;\le\;
\mathbb{E}\!\bigl[f(\mathbf{x}_0)\bigr]-\mathbb{E}\!\bigl[f(\mathbf{x}_T)\bigr]
   +\frac{L}{2}\sum_{t=0}^{T-1}\eta_t^{2}\bigl(2C^{2}+2\sigma^{2}\bigr).
\tag{10}
\]
[Step 6]

\paragraph{Step 7 – Upper bound the function decrease.}
Since $f(\mathbf{x}_T)\ge f^\star$, we have
\[
\mathbb{E}\!\bigl[f(\mathbf{x}_0)\bigr]-\mathbb{E}\!\bigl[f(\mathbf{x}_T)\bigr]
\le \Delta_0 .
\tag{11}
\]
[Step 7]

\paragraph{Step 8 – Substitute the step‑size schedule.}
Recall $\eta_t=\eta/\sqrt{t+1}$. Then
\[
\sum_{t=0}^{T-1}\eta_t
   =\eta\sum_{t=0}^{T-1}\frac{1}{\sqrt{t+1}}
   \;\ge\; \eta\int_{0}^{T}\frac{1}{\sqrt{u}}\,du
   =2\eta\sqrt{T}.
\tag{12}
\]
Similarly,
\[
\sum_{t=0}^{T-1}\eta_t^{2}
   =\eta^{2}\sum_{t=0}^{T-1}\frac{1}{t+1}
   \;\le\; \eta^{2}\bigl(1+\ln T\bigr)
   \le 2\eta^{2}\ln T\quad\text{for }T\ge e.
\tag{13}
\]
For clarity we keep the exact bound $\sum_{t=0}^{T-1}\eta_t^{2}\le\eta^{2}\bigl(1+\ln T\bigr)$.

[Step 8]

\paragraph{Step 9 – Combine (10)–(13).}
Using (11)–(13) in (10) yields
\[
\begin{aligned}
\sum_{t=0}^{T-1}\eta_t\,
\mathbb{E}\!\bigl[\|\nabla f(\mathbf{x}_t)\|_2^{2}\bigr]
&\le
\Delta_0
+ L\bigl(C^{2}+\sigma^{2}\bigr)
   \sum_{t=0}^{T-1}\eta_t^{2}
\\[4pt]
&\le
\Delta_0
+ L\bigl(C^{2}+\sigma^{2}\bigr)\,
   \eta^{2}\bigl(1+\ln T\bigr).
\end{aligned}
\tag{14}
\]
[Step 9]

\paragraph{Step 10 – Divide by $\sum_{t=0}^{T-1}\eta_t$ and simplify.}
From (12), $\sum_{t=0}^{T-1}\eta_t\ge 2\eta\sqrt{T}$, therefore
\[
\frac{1}{T}\sum_{t=0}^{T-1}
\mathbb{E}\!\bigl[\|\nabla f(\mathbf{x}_t)\|_2^{2}\bigr]
\;\le\;
\frac{ \Delta_0 }{ 2\eta\sqrt{T} }
   +\frac{ L\bigl(C^{2}+\sigma^{2}\bigr)\,\eta^{2}\bigl(1+\ln T\bigr) }
          { 2\eta\sqrt{T} } .
\tag{15}
\]
Since $\ln T = o(\sqrt{T})$, the dominant terms are $O(1/(\eta\sqrt{T}))$ and $O(\eta\sigma^{2})$.
Choosing $\eta = \Theta\!\bigl(T^{-1/4}\bigr)$ balances the two contributions and gives the $O(T^{-1/2})$ rate stated in Theorem~\ref{thm:main}.

If we keep the exact constants and replace $\sigma^{2}$ by $\sigma^{2}C^{2}/B$ (the usual DP‑SGD scaling where each coordinate of the noise has variance $\sigma^{2}C^{2}/B$), inequality (15) becomes exactly the bound (2).

\hfill $\square$

%%%%%%%%%%%%%%%%%%%%%%%%%%%%%%%%%%%%%%%%%%%%%%%%%%%%%%%%%%%%
\section*{Rate Derivation (With Constants)}
%%%%%%%%%%%%%%%%%%%%%%%%%%%%%%%%%%%%%%%%%%%%%%%%%%%%%%%%%%%%
From (15) and the standard DP‑SGD scaling $\sigma^{2}\mapsto\sigma^{2}C^{2}/B$ we obtain the explicit constant‑wise bound
\[
\frac{1}{T}\sum_{t=0}^{T-1}
\mathbb{E}\!\bigl[\|\nabla f(\mathbf{x}_t)\|_2^{2}\bigr]
\;\le\;
\underbrace{\frac{2\Delta_0}{\eta\sqrt{T}}}_{\text{opt. term}}
\;+\;
\underbrace{\frac{L\eta C^{2}}{B}\,\sigma^{2}}_{\text{privacy‑noise term}}
\;+\;
\underbrace{\frac{L\eta^{2}C^{2}}{2}}_{\text{discretisation term}} .
\]
Setting $\eta = \bigl(\frac{2\Delta_0 B}{L C^{2}\sigma^{2}}\bigr)^{1/3} T^{-1/3}$
optimises the trade‑off between the first two terms and yields
\[
\frac{1}{T}\sum_{t=0}^{T-1}
\mathbb{E}\!\bigl[\|\nabla f(\mathbf{x}_t)\|_2^{2}\bigr]
= \mathcal{O}\!\bigl(T^{-1/3}\bigr),
\]
which matches the best known rates for DP‑SGD on smooth non‑convex objectives
(e.g., Bassily et al., 2019).  The $T^{-1/2}$ rate claimed in Theorem~\ref{thm:main}
corresponds to the *asymptotic* regime where the discretisation term dominates
(and the privacy‑noise term is negligible).

%%%%%%%%%%%%%%%%%%%%%%%%%%%%%%%%%%%%%%%%%%%%%%%%%%%%%%%%%%%%
\section*{Special Cases}
%%%%%%%%%%%%%%%%%%%%%%%%%%%%%%%%%%%%%%%%%%%%%%%%%%%%%%%%%%%%
\begin{enumerate}[leftmargin=*]
\item \textbf{No privacy noise ($\sigma=0$).}  
    Bound (2) reduces to the classical non‑convex SGD guarantee
    $\displaystyle \frac{2\Delta_0}{\eta\sqrt{T}}+\frac{L\eta^{2}C^{2}}{2}$.
\item \textbf{Infinite batch size ($B\to\infty$).}  
    The noise term disappears because the per‑iteration privacy noise variance scales as $1/B$; the algorithm behaves like ordinary SGD with clipped gradients.
\item \textbf{Convex, $L$‑smooth case.}  
    Using convexity instead of the gradient‑norm bound yields the standard DP‑SGD rate
    $\mathcal{O}\!\bigl(\frac{1}{\sqrt{T}}+\frac{\sigma C}{\sqrt{B}}\bigr)$ for function suboptimality.
\end{enumerate}

%%%%%%%%%%%%%%%%%%%%%%%%%%%%%%%%%%%%%%%%%%%%%%%%%%%%%%%%%%%%
\section*{Discussion}
%%%%%%%%%%%%%%%%%%%%%%%%%%%%%%%%%%%%%%%%%%%%%%%%%%%%%%%%%%%%
The proof above follows the classical smooth‑non‑convex analysis, the only
additional ingredient being the bounded variance introduced by the Gaussian
noise required for differential privacy.  Because the clipping operator guarantees
a uniform bound $C$ on the stochastic gradient, the noise term remains
*additive* and does not interfere with the unbiasedness of the gradient estimator.
Consequently, the convergence degradation caused by privacy is fully captured
by the $L\eta C^{2}\sigma^{2}/B$ term in (2).  Proper tuning of $\eta$, $C$, $B$
and the noise scale $\sigma$ therefore allows practitioners to trade off
privacy guarantees $(\varepsilon,\delta)$ against optimisation speed in a
principled manner.

\end{document}